\subsection*{The marked genus-three source and an arithmetic class-two repair}

\begin{statement}{D.57}{Proved}{Both arithmetic branches have a full class-two repair over the actual order-3888 quotient}
Use the actual arithmetic element \(h\), the group
\(G_*=\mathbb F_3^4\rtimes Q_*\), and the maps of D.46,
D.49 and D.56. In particular \(|Q_*|=48\), \(|G_*|=3888\),
and the marked epimorphism sends \(s,r\) to the elements
specified in D.56. Put
\[
 K_Q=\ker(\widehat B\longrightarrow Q_*),\qquad
 U_Q=K_Q^{\mathrm{ab},(2)},\qquad
 L_{Q,2}=\overline{\Gamma_3K_Q(K_Q')^2K_Q^4},
 \qquad J_{Q,2}=\widehat B/L_{Q,2}.
 \tag{D.57.1}
\]
The subgroup \(K_Q\) has free rank nine. Its actual marked
compactification is
\[
 C:\ y^4=x^3+1,\qquad
 b(x,y)=\frac{(x^6+20x^3-8)^2}{64(x^3+1)^3}.
 \tag{D.57.2}
\]
The four points above infinity are removed. The ordered deck
generators are
\[
 s:(x,y)\longmapsto
       \left(\frac{2-x}{x+1},\frac{\sqrt3\,y}{x+1}\right),
 \qquad r:(x,y)\longmapsto(\zeta_3x,y).
 \tag{D.57.3}
\]
These formulas identify the actual 48-state marked quotient,
including both generator maps, with the deck group of (D.57.2).

Write \(D_t\) for the action of \(t=sr\) on \(U_Q\),
\(J=D_t^3\), and
\(N=\sum_{q\in[Q_*,Q_*]}q\). The derived group in this sum
is the marked quaternion group of order eight. The actual
arithmetic action of \(h\) on \(U_Q\) is one of the two
integral matrices
\[
 T_1=J+\frac{N/2-D_t^2N}{7},\qquad
 T_3=J^3+\frac{N/2-D_t^2N}{7}.
 \tag{D.57.4}
\]
Both matrices preserve the integral ribbon intersection form
over \(\mathbb Z_2\) and are invertible modulo two. Their
difference modulo two has rank four. The proof below treats
both matrices and the complete integral conjugator ambiguity;
no choice between them is required for the conclusion.

Let \(K_*\), \(M_*\), \(L_{*,2}\), \(U_*\), \(V_*\),
\(W_*\), and \(Z_*=W_*\wedge V_*\) be the actual objects of
D.56. There is an actual continuous peripheral automorphism
of
\[
 X_* =\widehat B/
       (\mathfrak R_h\cap M_*\cap L_{*,2})
 \tag{D.57.5}
\]
which fixes \(t\), has multiplier one, and induces exactly
D.49's prescribed automorphism on
\(\mathcal W_h=\widehat B/\mathfrak R_h\).
It induces the identity on \(G_*\) and an automorphism of the
full class-two exponent-four quotient
\(\mathcal T_{*,2}=\widehat B/L_{*,2}\).
Its action on the full integral homology quotient
\(\widehat B/M_*\) is a repaired action; the particular
homology representative chosen in D.50 is not asserted to
remain fixed. The exact underlying source and all projection
kernels are those of D.56.12. The resulting parameter belongs
to the image of \(\widehat F'\) and maps to the same literal
parameter \(f_h^*\) in \(\mathcal W_h\). Consequently the
previous metabelian and prescribed CM comparisons are retained
through those quotient maps.

Here are the evaluated dimensions in the actual arithmetic
source:
\[
\begin{array}{c|r}
 \dim V_*&649\\
 \dim W_*&640\\
 \dim(Z_*^{A_*})&2560\\
 \dim(Z_*)_{G_*}&74\\
 \operatorname{rank}(\text{full group norm})&40\\
 \operatorname{rank}(\text{compact group norm})&6\\
 \operatorname{rank}(\text{local same-homology correction})&64\\
 \dim(\text{reduced obstruction space})&10\\
 \dim(\text{compact-norm kernel/local image})&4 .
\end{array}
 \tag{D.57.6}
\]
The complete forty-variable equations retain all ten output
coordinates. In the recorded coordinates their coefficients
span a nine-dimensional subspace, with the explicit relation
that output coordinate zero equals output coordinate two.
The full lists are retained, rather than replacing the
ten-coordinate map by this smaller image.
For the two matrices in (D.57.4), the recorded zeros are
\[
\begin{array}{c|c|r|c}
 \text{branch}&\text{forty-parameter mask}&
       \text{fixed-seed trial}&\text{reduced constant}\\
 J&\mathtt{ba190a8d71}&682&\mathtt{040}\\
 J^3&\mathtt{dfd51862a6}&18&\mathtt{04d}.
\end{array}
 \tag{D.57.7}
\]
Each row has all eight integral \(s\)-conjugator framing cases,
with an explicit correction in the 600-dimensional local
parameter space. All sixteen corrected literal cusp words
have zero abelian error and zero in every one of the 74 moving
central coinvariant coordinates. Their lengths range from
114712 to 119132 letters in \(s,r,r^{-1}\). Every resulting
marked homology map has rank 649, preserves the intersection
pairing, and projects to its specified matrix in (D.57.4).

This proves the actual automorphism of (D.57.5). The next
full-source map still has the explicit kernel
\[
 \widehat B\longrightarrow X_*:
       \mathfrak R_h\cap M_*\cap L_{*,2}.
 \tag{D.57.8}
\]
No surjectivity from genuine full profinite \(GT_0\) is inferred.
The exact map from the original full \(G_3\) class-two source
of D.54 to (D.57.5) is induced by
\[
 \mathfrak R_h\cap M_\infty\cap L_2
 \ \subset\
 \mathfrak R_h\cap M_*\cap L_{*,2}.
 \tag{D.57.9}
\]
Its kernel is the quotient of the right-hand subgroup by
the left-hand subgroup. The original full \(G_3\) invariant
equation remains to be lifted across this map.

The manuscript passages are Part D(e), section 51.16,
folio 734 verso, scan 174, and the two-torus/Fermat
follow-ups in section 51.15, folios 725--727, scans 156--160.
\end{statement}
\begin{proof}
\emph{The exact marked curve.}
Use homogeneous coordinates \((X,Z)\), with
\(Y^4=X^3Z+Z^4\). The matrices for the two maps in
(D.57.3) are
\[
 S=\frac1{\sqrt3}
     \begin{pmatrix}-1&2\\1&1\end{pmatrix},
 \qquad R=\begin{pmatrix}\zeta_3&0\\0&1\end{pmatrix},
 \tag{D.57.10}
\]
and \(Y\) is unchanged. Both matrices preserve the displayed
quartic, \(S^2=R^3=I\), and \((SR)^3=iI\).
The latter scalar acts on the affine curve by
\((x,y)\mapsto(x,-iy)\), so it is exactly the operator \(J\).
Direct substitution shows that \(b\) is invariant under
both maps. The useful second identity is
\[
 b-1=\frac{x^3(x^3-8)^3}{64(x^3+1)^3}.
 \tag{D.57.11}
\]
The branch permutations are therefore \(2^{24}\),
\(3^{16}\), and \(12^4\) above zero, one, and infinity.
Their degree is 48. The exact generated matrix group also
has 48 elements. The program constructs its bijection with
the 48 matrices of D.56 and verifies all 96 ordered
\(s,r\) edges in both directions. Thus this comparison
identifies the marked epimorphism, not only the abstract
group order. On the positive real interval, the \(s\)-fixed
point has \(x=\sqrt3-1\), while the \(r\)-fixed point has
\(x=0\); the local multipliers are respectively \(-1\)
and \(\zeta_3\). These fix the branch convention in (D.57.3).
The genus is three and there are four cusps, giving free
rank \(2\cdot3+4-1=9\).

The quotient by \(y\mapsto-y\) is the elliptic curve
\(E_3:Y^2=x^3+1\), through \(Y=y^2\). The natural
index-eight map to the elliptic source of D.45 is this
map followed by multiplication by two. Explicitly it is
\[
 u=\frac{x(x^3-8)}{4(x^3+1)},\qquad
 v=\frac{x^6+20x^3-8}{8y^6},\qquad
 v^2=u^3+1,\quad b=v^2.
 \tag{D.57.12}
\]
Substitution gives \(s:(u,v)\mapsto(u,-v)\) and
\(r:(u,v)\mapsto(\zeta_3u,v)\). Its deck kernel is
therefore exactly the marked quaternion group, since
these images generate the abelian quotient of order six.
In particular \(t^2\) acts by \(u\mapsto\zeta_3^2u\),
which proves the precise CM evaluation convention used
in (D.57.4). This also identifies the rational Tate map
with D.46's elliptic source; no unidentified elliptic
twist is introduced. The two further degree-two
quotient maps can also be written explicitly.
For \(a=\sqrt3\) or \(-\sqrt3\), put
\[
 X_a=\frac{x+1-a}{x+1+a},\qquad
 u_a=\frac{y}{x+1+a},\qquad
 z_a=\frac{(9+6a)X_a^2-9}{6a}.
\]
Then
\[
 z_a^2=1-(12+8a)u_a^4,\qquad
 E_a:\ v^2=w^3+(3+2a)w,
 \tag{D.57.13}
\]
where the rational map to \(E_a\) is obtained first by
\(w_0=2(1+z_a)/u_a^2\),
\(v_0=4(1+z_a)/u_a^3\), then by
\(w=w_0/4\), \(v=v_0/8\).
All identities are verified in the function field of (D.57.2).
The pulled-back differentials for the two quartics are
\[
 \frac{du_a}{z_a}=
       \frac{-x-1+a}{4}\frac{dx}{y^3}.
 \tag{D.57.14}
\]
Together with \(dx/y^2\) they are independent. The induced
homomorphism of Jacobians is consequently an isogeny
\(\operatorname{Jac}(C)\to E_3\times E_{\sqrt3}\times
E_{-\sqrt3}\): its differential is an isomorphism, hence
its image has full dimension and its kernel is finite.
Each \(E_a\) becomes isomorphic over \(\overline{\mathbb Q}\)
to \(E_i:v^2=w^3-w\) by scaling by a fourth root of
\(-(3+2a)\). This gives the actual comparison of the
four-dimensional complementary Tate module with two copies
of the Tate module of \(E_i\).

\emph{Why there are exactly two integral arithmetic branches.}
On the nine-dimensional rational homology the simultaneous
centralizer of the two marked deck matrices has dimension
five. Here this is an exact calculation: the 81-variable
commutator matrix has rank 76 modulo five, and the five
independent rational operators
\(I,J,J^2,J^3,D_t^2N\) commute with both deck matrices.
The three mutually orthogonal projectors are
\[
 e_3=N/8,\qquad e_i=(I-J^2)/2,\qquad
 e_c=(I+J^2)/2-N/8,
 \tag{D.57.15}
\]
of ranks two, four, and three. The last is the cusp subspace;
the first is the marked \(E_3\) subspace. On the middle
subspace the centralizer is exactly \(\mathbb Q_2(i)\),
with \(i\) represented by \(J\).

There is an actual finite-action argument for the middle
arithmetic scalar. The curve \(E_i\) has good reduction at
three: its displayed integral model has discriminant 64.
Its \(2^n\)-torsion is finite etale at three for every \(n\),
so the inertia element \(\tau\) used in D.46 acts trivially
on \(T_2E_i\). Therefore the actual commutator
\(h=[g,\tau]\) also acts trivially there. The isogeny just
constructed supplies the evaluation map
\[
 \operatorname{Hom}_{\overline{\mathbb Q}}
       (E_i,\operatorname{Jac}(C)^-)
       \otimes_{\mathbb Z[i]}T_2E_i
       \longrightarrow T_2\operatorname{Jac}(C)^-
 \tag{D.57.16}
\]
as an isomorphism after tensoring with \(\mathbb Q_2\).
Its Hom lattice is finitely generated: under a complex
embedding it embeds in the integral homomorphisms between
the first homology lattices, because an abelian-variety
homomorphism is determined by its analytic lattice map.
Its Galois action is therefore finite: choose a finite
integral generating set of morphisms, each defined over a
finite extension, and take their common field of definition.
All deck maps are defined over \(\mathbb Q(\zeta_{12})\).
The commutator \(h\) fixes this field. Passing from the
arithmetic action on the curve to the marked source action
introduces a deck element; since the latter action is the
identity on \(Q_*\), that deck element is central.
Equivalently the finite Galois image on the Hom lattice
normalizes the finite deck image, so their combined image
is finite. It follows from (D.57.16) that the marked scalar
on the middle summand has finite order. The only roots of
unity in \(\mathbb Q_2(i)\) are \(1,-1,i,-i\): the
residue field has order two, and its ramified quadratic
extension contains \(\mu_4\) but not \(\mu_8\);
indeed \(\Phi_8(X+1)\) is Eisenstein at two and has degree four.

On the first summand, D.46 gives the actual marked scalar
\(\eta_*=(-3-8\zeta_3^2)/7\), with \(\zeta_3^2\)
represented by \(D_t^2\). On the cusp summand the action
is the identity, because the source action fixes the
marked deck quotient and has multiplier one. Consequently
the four rational candidates are exactly
\[
 J^k+\left(\frac{-3I-8D_t^2}{7}-(-1)^kI\right)N/8,
       \qquad 0\leq k<4.
 \tag{D.57.17}
\]
The marked integral graph lattice excludes \(k=0,2\).
For both of those candidates the entry in row one and
column zero, with zero-based indices in the certificate,
is \(-1/14\). The other two candidates have integral
\(\mathbb Z_2\) entries and are precisely (D.57.4).
The certificate checks their exact rational pairing
identities, ranks modulo two, and all marked deck
commutations. This argument uses the two-adic Hom and Tate
comparison (D.57.16); it does not identify a two-adic scalar
from the ramified three-adic representation.

\emph{The actual common quotient and the remaining framing.}
The exact sequence
\(1\to K_*\to K_Q\to A_*\to1\), with \(|A_*|=81\),
gives the integral onto map \(U_*\to U_Q\).
Odd averaging identifies its kernel with \(U_*^-\).
Thus, already on full homology, the marked common quotient
is
\[
 (\widehat B/M_*)/U_*^-
       \simeq G_*\mathbin{\times}_{Q_*}J_{Q,1},
 \qquad J_{Q,1}=\widehat B/\ker(K_Q\to U_Q).
 \tag{D.57.18}
\]
Both projections are the original source projections.
For the class-two quotient let \(D_*=\ker(\mathcal T_{*,2}
\to G_*)\) and let \(E_*\) be the image of
\(\mathfrak R_h\) in \(\mathcal T_{*,2}\).
D.51's averaged quotient construction, applied through
D.56's actual maps, gives
\(E_*=[D_*,\widetilde A_*]\),
\(E_*\cap D_*'=Z_*\), and
\[
 \mathcal T_{*,2}/E_*
        \simeq G_*\mathbin{\times}_{Q_*}J_{Q,2}.
 \tag{D.57.19}
\]
Here are both directions of this identification. In the
preimage of \(A_*\), an odd complement exists and the
maximal two-group quotient kills it and its commutators
with \(D_*\). The resulting group is the free class-two
exponent-four quotient of \(K_Q^{(2)}\), of rank nine.
Conversely its defining laws already hold on the image
of \(K_*\), so its quotient map factors through
\(\mathcal T_{*,2}\). Its kernel on \(D_*\) is exactly
\(E_*\), by the averaged quotient and commutator calculation
of D.51. The two remaining kernels over \(Q_*\) have
orders a power of two and a power of three; finite Goursat
therefore gives the full displayed fibre product.

The actual \(\mathcal W_h\) action on (D.57.19) is
\((\mathrm{id},\alpha_h)\). The first projection factors
through the identity three-primary component. For the
second, \(C_h\) maps into the two-group kernel over
\(Q_*\), so its map factors through the two-primary
component, where the action is exactly \(h\).
This proves the common action on the full class-two
quotient, before using its homology matrices.

There is also an exact way to remove the homology-level
based ambiguity without discarding it. The integral graph
sequence gives
\(U_Q^{Q_*}=0\) and \(H^1(Q_*,U_Q)=0\).
For the latter assertion the invariant vertex kernel is
generated by \((2,-3)\), since the vertex orbit sizes
are 24 and 16; this is exactly the boundary of the
invariant edge norm. The two vertex permutation lattices
and the edge lattice have zero first cohomology, as their
stabilizers are finite and the coefficients are torsion
free. Hence two automorphisms of \(J_{Q,1}\) with the
same action on \(U_Q\) and the same identity action on
\(Q_*\) differ by a unique inner element of \(U_Q\).
When both fix \(t\), that element is in \(U_Q^t\).
The exact marked graph calculation gives
\[
                 U_Q^t=\mathbb Z_2[t^{12}],
 \tag{D.57.20}
\]
including primitivity of the displayed cusp class.
Thus the comparison is the actual map
\(\alpha\mapsto\operatorname{Int}(t^{-12z})\alpha\)
for one \(z\in\mathbb Z_2\). This inner automorphism
exists in every source under discussion and fixes \(t\).
We construct the repair after this comparison, then
compose with \(\operatorname{Int}(t^{12z})\) to recover
the original prescribed \(\mathcal W_h\) action.

After this exact comparison, conjugators for the same
\(J_{Q,1}\) action differ by the fixed lattices
\(U_Q^s\) and \(U_Q^r\). The Smith decompositions of
\(s-1\) and \(r-1\) are respectively
\[
 (1,1,1,2,2,2,0,0,0),\qquad
 (1,1,1,1,1,3,3,3,0).
 \tag{D.57.21}
\]
Their unimodular change-of-basis matrices are retained;
the resulting saturated integral fixed ranks are three
and one. In particular these are the integral fixed
lattices, not the larger fixed spaces obtained by
reducing the two-primary action modulo two.

The finite \(s\)-conjugator coordinate in \(Q_*\) is
one of the two order-four elements 26 and 33 recorded
in D.56's ordered matrix list. Indeed a normalized
conjugator lies in \(F\), centralizes \(s\), and its
finite CM coordinate excludes the central elements.
The centralizer in \(Q_*\) is a two-group, so its even
part lies in the quaternion derived group. Exact
integral equations pair state 26 with \(T_1\), and
state 33 with \(T_3\). The other two pairings have
no integral solution. The normalized \(r\)-conjugator
has quotient coordinate one, by D.49's coprime
centralizer correction. Above each of the two
\(s\)-coordinates there are nine elements of
\(C_{G_*}(s)\). Any prescribed one can be used:
the actual joint quotient of the three-primary factor
with \(\mathcal H_h\) is their fibre product over
\(Q_*\), and coprime order-two fixed-coset lifting
through its pro-three kernel lifts the chosen
centralizing pair. This only changes a conjugator
within its centralizer and leaves the prescribed
automorphism and parameter unchanged.

\emph{The complete arithmetic central module.}
All calculations now use the 3888-state marked graph.
There are 649 explicit free words. The 40 nonzero
character pairs of \(A_*\simeq\mathbb F_3^4\) split
\(W_*\) into 40 summands of dimension 16 over
\(\mathbb F_2\). For a character \(\ell\), the
projector is \(e_{\ker\ell}+e_{A_*}\); a generator
of \(A_*/\ker\ell\) acts by
\((u,v)\mapsto(v,u+v)\) on each of eight pairs.
The invariant exterior square has the following
64-element basis:
\[
 u_i\wedge v_i,\quad
 u_i\wedge u_j+u_i\wedge v_j+v_i\wedge v_j,\quad
 u_i\wedge v_j+v_i\wedge u_j\quad(i<j).
 \tag{D.57.22}
\]
Cross terms between different character pairs and
terms with \(V_*^{A_*}\) have zero \(A_*\)-invariants.
Thus the actual invariant module has dimension
\(40\cdot64=2560\). The program verifies every
projection column and both marked generator actions,
then reduces the columns of \(1-s\) and \(1-r\).
Their rank is 2486, giving the actual 74-dimensional
coinvariant module. It also evaluates all 48 terms
of the remaining deck norm on every coinvariant
representative. The full norm rank is 40, while
compactification lowers the rank to six. This gives
all norm dimensions in (D.57.6) by actual maps.

Write
\(L(b,a)=(1-s)b+s(1-r)a\) on \(W_*\oplus W_*\).
It is onto. Its kernel \(P_*\) has dimension 640,
and the exact same-homology kernel is
\[
 K_*^{\mathrm{par}}=
 \{(u_s+d,u_r+d):
           u_s\in W_*^s,\ u_r\in W_*^r,\ d\in W_*^t\}.
 \tag{D.57.23}
\]
The three dimensions are 336, 210, and 54; their
displayed injection gives total dimension 600.
An exact complementary basis of dimension 40 in
\(P_*\) is retained in the certificate.
Each parameter is lifted into the integral moving
summand modulo four by the actual 81-term averaging
operator. Its even cusp error is removed with the
onto map \(L\), yielding an actual integral
peripheral pair. This verifies the entire reduction
map rather than assigning formal mod-two parameters.

The class-two collection model uses coordinates
\((a,c)\), with \(a\in(\mathbb Z/4)^{649}\) and
\(c\in\bigwedge^2\mathbb F_2^{649}\). Appending
generator \(j\) adds the coefficients
\(a_i\) with \(i>j\) to its central coordinates.
The program pulls back all 74 coinvariant functionals
through (D.57.22), and constructs the exact
\(s,t\) actions on this collected group from the
649 literal conjugated words. Its resulting local
map has rank 64 for each arithmetic seed. Every
local column is killed by the compact norm.
Since that norm has kernel dimension 68, the
four-dimensional residual in (D.57.6) is retained.

\emph{The forty-variable equations and their literal zeros.}
For a fixed homology seed, let \(d_s\) be the
abelian coordinate of its image of \(s\), multiplied
by \(s^{-1}\). If \(q_{\mathrm{id}}(b,a)\) is the
collected cusp error for the identity seed, the
actual affine equation is
\[
 q_0+q_{\mathrm{id}}(b,a)
       +[d_s\wedge s(a-b)]
       \quad\hbox{in }(Z_*)_{G_*}.
 \tag{D.57.24}
\]
To verify the formula, the corrected cusp product
in the kernel is
\(b\,s_0(b^{-1}a)\,t(a^{-1})\), multiplied by
the original central error. On central coordinates
the action of \(s_0\) differs from that of \(s\)
by exactly the commutator with \(d_s\).
The collected-product polar term is
\[
       [(1+s)b'\wedge s(a+b)] .
 \tag{D.57.25}
\]
Both formulas are evaluated on the actual integral
pair lifts. In particular changing central
representatives adds only augmentation columns
and cannot change (D.57.24).

The exact local image in (D.57.24) has rank 64.
Quotienting by this image gives the stated ten
coordinates. The certificate contains the
74-to-10 quotient matrix, all 40 linear
coefficients, all 780 quadratic coefficients,
and the constant, separately for both branches.
The saturated \(s\)-fixed lattice of (D.57.21)
has three generators. Each has zero obstruction
in all 74 coordinates, as does its double.
The integral \(r\)-fixed generator and its
double also have zero obstruction. Cross terms
with moving parameters vanish: the plus and
moving summands pair to zero under the odd
norm, and (D.57.25) gives the same conclusion
directly. Thus its complete reduction consists
of the eight explicitly retained \(s\)-framing
cases, with no unrecorded higher two-adic choice.

The search evaluates the full ten-coordinate
polynomial and stops at the actual trials in
(D.57.7). These are existence searches, not an
assertion that all \(2^{40}\) inputs were tested.
For each of the sixteen cases, Gaussian elimination
then supplies the actual 600-parameter local
correction. The corrected words are evaluated
anew in the full collected group. They have
zero abelian error and zero under all 74
functionals. Independently the literal edge
cycle map is evaluated on all 649 generators;
it has full rank, preserves the actual ribbon
form, commutes with the two marked deck maps,
and has the precise rank-nine projection
(D.57.4). These are the tests reported in the
statement, including their recorded word lengths.

\emph{Alignment to the actual arithmetic class-two action.}
It remains to check that these numerical zeros
lift the prescribed source action. The common
quotient is the full quotient (D.57.19), on
which the actual arithmetic action already
exists. Its conjugators have one of the two
rank-nine matrices and one of the saturated
integral framing classes just exhausted.
Choose their actual lifts to \(\mathcal T_{*,2}\).
The onto moving homology map permits their
abelian coordinates to be the corresponding
ones in the certificate. If the selected
literal words differ in their plus central
coordinates, correct this difference by the
actual onto map
\[
 \bigwedge^2 V_*\longrightarrow
          \bigwedge^2 V_*^{A_*}.
 \tag{D.57.26}
\]
It is the commutator map of the averaged
quotient in D.51, not a new arbitrary quotient.
All 74 functionals kill this plus summand,
and changes of central conjugators give only
augmentation columns. Hence the computed
74-coordinate zero is unchanged.

After this alignment the cusp error maps to
one in (D.57.19), since the prescribed action
there fixes \(t\). It is central with zero
abelianization, so it belongs to
\(E_*\cap D_*'=Z_*\). Its zero image in
\((Z_*)_{G_*}\) says exactly that it lies
in the augmentation image
\((1-s)Z_*+s(1-r)Z_*\). Choose those two
central corrections and add them to the
conjugators. The cusp identity is now exact
in \(\mathcal T_{*,2}\). The two finite-order
relations are preserved because the images
remain conjugates of \(s\) and \(r\).
The rank-649 calculation proves generation
of the finite two-group kernel by Frattini;
the map on \(G_*\) is the identity, so this
is an automorphism of the whole finite quotient.

The same correction continues through the
full integral homology tower. The moving
integral cusp map is onto by odd averaging
and the augmentation formula for \(L\).
At stage \(2^k\), with \(k\geq2\), correct
the next cusp error by \(2^k\) times a
moving integral pair. Such corrections can
be chosen as fourth powers of actual
\(\mathfrak R_h\) lifts, and therefore
are one in \(\mathcal T_{*,2}\). They retain
the already constructed finite class-two
action and the entire \(\mathcal W_h\)
action. Compactness gives compatible
conjugators at all levels. The map is
invertible on the full two-adic homology
kernel because its reduction modulo two
is invertible. The exact fibre-product
description of D.56 therefore gives the
claimed continuous automorphism of (D.57.5).
Finally undo the explicitly recorded
cusp-inner comparison (D.57.20).

After the initial common-quotient alignment,
all changes used to repair the prescribed
action lift in \(\mathfrak R_h\subset\widehat F'\).
Thus the final normalized parameter still
belongs to the image of \(\widehat F'\)
and maps exactly to \(f_h^*\). This proves
the retained metabelian and CM statements
through their existing source maps.

For clarity, the remaining source kernel
in (D.57.8) is nontrivial. Choose an actual
element \(u\in\mathfrak R_h\) whose image
in \(U_*^-\) is a primitive moving vector,
using D.50/D.56's onto map, and a free
generator \(v\in K_*\) independent of it
modulo two. Then \([u,v]^2\) belongs to
\(\mathfrak R_h\cap M_*\cap L_{*,2}\).
In the free class-two quotient with
abelianization \((\mathbb Z/8)^{649}\)
and commutator module
\(\bigwedge^2(\mathbb Z/4)^{649}\), its
image is \(2\bar u\wedge\bar v\ne0\).
This supplies an actual surviving source
element and the explicit comparison detecting
it. The inclusion (D.57.9) follows from
\(M_\infty\subset M_*\) and functoriality
of the defining class-two exponent-four
laws for \(K_3\subset K_*\). It records
the next unsolved lifting map without
asserting a full profinite realization.
\end{proof}
